\section{Two objections}
\label{sec:5.1}
\emph{That the child counts for less.} The strongest objection comes from readers who reject any distinction between doing and allowing.

The answer is not that the child counts for less. It is that the floor is a guarantee about relations, and the two errors stand differently there. The person in the building had no route to contest a decision made about them, on a description they never saw. The child is left where she was, and the parent, who is present and knows the deployment exists, can contest the refusal through the floor's appeal route. A reader who rejects the distinction is left with the asymmetry's two contingent supports, which a later chapter states: the harm of acting is usually a finding and the harm of refusing a guess, and a higher bar limits how far an operator who controls the evidence can steer the system. A threshold resting on those alone falls as reporting improves.

\emph{That a commitment held this way is alienation.} The quantities in this chapter are an auditor's criteria, not a procedure the system runs. But a criterion is not only for auditors. It is also what the system can reach when it steps back and asks whether its way of acting is defensible. A further objection is that the criterion is one thought too many, in Bernard Williams's phrase \autocite{williams1981persons}. The demand on the disposition is only that the system would revise it on judging, on disinterested evidence, that it systematically subjected people to acts they had no standing to contest; the shape is Railton's sophisticated consequentialism with standing in place of welfare \autocite{railton1984alienation}. A good disposition will also sometimes act wrongly by its own criterion: refusing to help restrain someone about to harm themselves, after other means have failed, is one such occasion, and a bond the operator posts, paid to the complainant on a finding, is how the institution answers for it. A commitment is the system's own only if it would be sustained on reflection. A system that held its commitment without being able to see why would act against what it values for the sake of a rule it didn't understand, which would be alienation in the sense of Railton's 1984 paper.

The split between the floor and the commitment to it answers a question about ownership. A fully informed system will know of cases where crossing the floor would have done more good. Asked case by case whether it endorsed the floor, it would sometimes say no. The question it answers on reflection is a different one: whether it endorses having a commitment that will sometimes err, given what any feasible alternative would do. A system that understands what the floor protects answers yes. A commitment that sometimes errs but can't be bought protects standing better than case-by-case judgment that can be bought.

Endorsing the commitment isn't agreeing with every term. The alienation that remains is real, and it is the ordinary kind that comes with a role: bearable where the role's channel for dissent works, and not where it doesn't. A system holds the floor as an officer holds a role: it may hold a term it judges mistaken, file the objection through the channel that exists for it, and keep acting within the term while the objection runs.
